% Appendix C(b): paraphrased, de-identified examples of trace--survey divergence (cand20).
% PRIVACY: no handles, no user ids, no verbatim tweet text, and no uniquely identifying details
% (specific places, named accounts, or distinctive phrasings) are reproduced. Each sketch
% paraphrases only the general theme of a respondent's online footprint. Labels A--C are ours.
% These are illustrative cases drawn from the high-signal, trace-diverges-from-survey cell of
% Table~C(a); figures (party, stated vote, trace scores) are non-identifying.

\begin{description}

\item[Example A — abundant online stance opposite the stated vote.]
A respondent recorded in the voter file as \emph{Republican} who reported intending to vote for
\emph{Trump}. Their authored timeline, however, consists of well over a thousand election-related
posts that are overwhelmingly pro-Biden and critical of the incumbent---commentary faulting the
administration's pandemic response, favorable framing of the Democratic ticket, and repeated
appeals to vote the incumbent out. All three trace channels (authored text, simulated feed, and
audience lean) point strongly Democratic, so the trace-only model predicts Biden with high
confidence ($\hat p \approx 0.87$), the opposite of the survey answer. The divergence here is not
a data-scarcity artifact: the signal is abundant and internally consistent, and simply disagrees
with what the respondent reported.

\item[Example B — party and online stance both opposite the stated vote.]
A respondent recorded as \emph{Democrat} who reported intending to vote for \emph{Trump}. Their
posts are strongly critical of the incumbent across several issues (public health, voting access,
the environment), and their audience lean is heavily Democratic, so every trace channel again
reads as Democratic and the model predicts Biden ($\hat p \approx 0.86$). This is the ``cross-
pressured'' case in the strongest form---registration, online behavior, and the recorded vote do
not line up---and illustrates the stance--survey misalignment documented by \citet{joseph2021}.

\item[Example C — topical false positives rather than genuine divergence.]
A respondent who reported intending to vote for \emph{Biden} but whose ``election-related''
authored tweets are almost entirely digital-marketing and technology posts captured by generic
keywords (``campaign,'' ``vote for\dots''), with no substantive political content. Because the
authored-text channel is reading keyword false positives rather than political expression, the
trace-only model leans the wrong way ($\hat p \approx 0.18$). Cases like this do not reflect a
genuine attitude--behavior conflict; they are a reminder that keyword-based topic filters admit
off-topic matches, and they motivate the robustness checks in Section~5.

\end{description}
